\chapter{正交与最小二乘法}

在这一章中，我们将重点分析$m>n$的高矩阵$\vb{A}\in\R^{m\times n}$的分解问题。该问题适用本章引入的QR分解$\vb{A}=\vb{Q}\vb{R}$，其中$\vb{Q}\in\R^{m\times m}$是正交矩阵而$\vb{R}\in^{m\times n}$是高的上三角矩阵。首先，我们将介绍两种重要的正交变换：代表反射的Householder矩阵，代表旋转的Givens矩阵。随后将引入QR分解的概念并证明其存在性和唯一性，接下来将会探讨基于Householder变换和Givens变换的QR分解算法，以及经典的基于Gram-Schmidt变换的QR分解算法。最后我们会探讨QR分解的应用：在$\vb{A}$超定的情况下解决$\min_{\vb{x\in\R^n}}\norm{\vb{y}-\vb{A}\vb{x}}_2^2$的最小二乘问题。

\section{两种正交变换}

我们知道，正交矩阵$\vb{Q}\in\R^{m\times m}$是指满足以下条件的矩阵
\begin{Equation}[正交矩阵定义]
    \vb{Q}^T\vb{Q}=\vb{Q}\vb{Q}^T=\vb{I}
\end{Equation}

在LU分解中，消元依靠的是高斯变换。在QR分解中，消元依靠的则是正交变换！可以这样理解：QR分解同样要产生上三角矩阵$\vb{R}$，该过程就是由一系列正交变换构成的$\vb{Q}$实现的。

在$\R^{2\times 2}$下，有两种重要的正交变换：反射和旋转。对于$\theta$，我们记
\begin{Equation}[2x2正交cs]
    c=\cos\theta\qquad s=\sin\theta
\end{Equation}

旋转（Rotation）矩阵$\vb{Q}$可以表示为以下形式，代表旋转$\theta$的角度
\begin{Equation}[2x2正交旋转]
    \vb{Q}=\mqty[
        c & s \\
        -s & c \\
    ]
\end{Equation}

反射（Reflection）矩阵$\vb{Q}$可以表示为以下形式，代表沿着$\theta/2$反射
\begin{Equation}[2x2正交反射]
    \vb{Q}=\mqty[
        c & s \\
        s & -c \\
    ]
\end{Equation}

这一节要研究的问题是，如何将$\R^{2\times 2}$上的旋转和反射以某种方式应用到$\R^{m\times m}$的矩阵上。

\subsection{Householder反射变换}

在$\R^m$下，通常反射会选取一个法向量$\vb{v}\in\R^m$并讨论$\vb{x}\in\R^m$对于垂直$\vb{v}$的超平面的镜像。

\begin{BoxFormula}[Householder反射变换]
    Householder反射变换矩阵$\vb{H}\in\R^{m\times m}$可以表示为以下形式
    \begin{Equation}[Householder反射变换]
        \vb{H}(\vb{v})=\vb{I}-\beta\vb{v}\vb{v}^T\qquad\beta=\frac{2}{\vb{v}^T\vb{v}}
    \end{Equation}
    其中$\vb{v}\in\R^m$为反射超平面的法向量。
\end{BoxFormula}

现在的问题是，对于任意的$\vb{x}$，如何找到对应的$\vb{v}$构造出变换矩阵$\vb{H}$，使$\vb{H}\vb{x}=\vb{e}_1$成立，换言之，这就是将$\vb{x}$变换为只有第一个元素为$1$其余元素都为$0$的单位向量$\vb{e}_1$，即所谓消元！
\begin{Equation}[反射变换1]
    \mqty[
        \vb{x}(1)\\
        \vb{x}(2)\\
        \vdots\\
        \vb{x}(m)\\
    ]\to\mqty[
        \times\\
        0\\
        \vdots\\
        0
    ]
\end{Equation}

根据\cref{fml:Householder反射变换}可得，请注意由于$\vb{v}^T\vb{x}$是数，所以可以提到前面
\begin{Equation}[反射变换2]
    \vb{H}\vb{x}=\qty(\vb{I}-\frac{2\vb{v}\vb{v}^T}{\vb{v}^T\vb{v}})\vb{x}=\vb{x}-\frac{2\vb{v}^T\vb{x}}{\vb{v}^T\vb{v}}\vb{v}
\end{Equation}

由于$\vb{v},\vb{x},\vb{e}_1$的反射关系，三者共平面，故可以将$\vb{v}$写为$\vb{x},\vb{e}_1$的线性组合
\begin{Equation}[反射变换3]
    \vb{v}=\vb{x}+\alpha\vb{e}_1
\end{Equation}

注意到\cref{eq:反射变换2}的分子和分母有$\vb{v}^T\vb{x}$和$\vb{v}^T\vb{v}$，现在要用\cref{eq:反射变换3}表示这两者。

项$\vb{v}^T\vb{x}$可以表示为
\begin{Equation}[反射变换4]
    \vb{v}^T\vb{x}=\vb{x}^T\vb{x}+\alpha\vb{x}(1)
\end{Equation}

项$\vb{v}^T\vb{v}$可以表示为
\begin{Equation}[反射变换5]
    \vb{v}^T\vb{v}=\vb{x}^T\vb{x}+2\alpha\vb{x}(1)+\alpha^2
\end{Equation}

现在将\cref{eq:反射变换3,eq:反射变换4,eq:反射变换5}代入\cref{eq:反射变换2}
\begin{Equation}[反射变换6]
    \vb{H}\vb{x}=1-2\frac{\vb{x}^T\vb{x}+\alpha\vb{x}(1)}{\vb{x}^T\vb{x}+2\alpha\vb{x}(1)+\alpha^2}\qty(\vb{x}+\alpha\vb{e}_1)
\end{Equation}

稍做整理
\begin{Equation}[反射变换7]
    \vb{H}\vb{x}=\qty(1-2\frac{\vb{x}^T\vb{x}+\alpha\vb{x}(1)}{\vb{x}^T\vb{x}+2\alpha\vb{x}(1)+\alpha^2})\vb{x}-2\alpha\frac{\vb{v}^T\vb{v}}{\vb{v}^T\vb{v}}\vb{e}_1
\end{Equation}

化简
\begin{Equation}[反射变换8]
    \vb{H}\vb{x}=\qty(\frac{\alpha^2-\vb{x}^T\vb{x}}{\vb{x}^T\vb{x}+2\alpha\vb{x}(1)+\alpha^2})\vb{x}-2\alpha\frac{\vb{v}^T\vb{v}}{\vb{v}^T\vb{v}}\vb{e}_1
\end{Equation}

请注意$\vb{v}^T\vb{v}$实际上就是$\norm{\vb{v}}_2^2$
\begin{Equation}[反射变换9]
    \vb{H}\vb{x}=\qty(\frac{\alpha^2-\norm{\vb{v}}_2^2}{\vb{x}^T\vb{x}+2\alpha\vb{x}(1)+\alpha^2})\vb{x}-2\alpha\frac{\vb{v}^T\vb{v}}{\vb{v}^T\vb{v}}\vb{e}_1
\end{Equation}

我们注意到，如果此时令$\alpha=\pm\norm{\vb{x}}_2$，就有$\vb{v}=\vb{x}\pm\norm{\vb{x}}_2\vb{e}$以及
\begin{Equation}[反射变换10]
    \vb{H}\vb{x}=-2\alpha\frac{\vb{v}^T\vb{v}}{\vb{v}^T\vb{v}}\vb{e}_1
\end{Equation}

现在已经知道$\vb{H}\vb{x}$平行于$\vb{e}_1$了，但到底$\vb{H}\vb{x}$是多少倍的$\vb{e}_1$呢？代入\cref{eq:反射变换4,eq:反射变换5}
\begin{Equation}[反射变换11]
    \vb{H}\vb{x}=-2\alpha\frac{\vb{x}^T\vb{x}+\alpha\vb{x}(1)}{\vb{x}^T\vb{x}+2\alpha\vb{x}(1)+\alpha^2}\vb{e}_1
\end{Equation}

代入$\alpha=\pm\norm{\vb{x}}_2$和$\vb{x}^T\vb{x}=\norm{\vb{x}}_2^2$
\begin{Equation}[反射变换12]
    \vb{H}\vb{x}=\mp\norm{\vb{x}}_2\vb{e}_1
\end{Equation}

我们将这个结论整理如下
\begin{BoxFormula}[Householder矩阵的构造]
    对于$\vb{x}\in\R^m$，若选取以下Householder向量$\vb{v}\in\R^m$
    \begin{Equation}[Householder矩阵的构造1]
        \vb{v}=\vb{x}\pm\norm{\vb{x}}_2\vb{e}_1
    \end{Equation}
    则$\vb{v}$构造的Householder矩阵满足
    \begin{Equation}[Householder矩阵的构造2]
        \vb{H}\vb{x}=\mp\norm{\vb{x}}_2\vb{e}_1
    \end{Equation}
\end{BoxFormula}

如\cref{fig:Householder反射变换}所示，请注意$\vb{x}$是关于垂直于$\vb{v}$的超平面（而不是关于$\vb{v}$）进行镜像的。

\begin{Figure}[Householder反射变换]
    \includegraphics[scale=0.9]{HouseholderReflect.fig.pdf}
\end{Figure}

通常来说我们会选取$\vb{v}=\vb{x}-\norm{\vb{x}}_2\vb{e}_1$以便得到$\vb{H}\vb{x}=\norm{\vb{x}}_2\vb{e}_1$的正的变换结果。

\begin{BoxAlgorithm}[Householder参数]
    对于$\vb{x}\in\R^m$，该函数计算满足$\vb{v}(1)=1$的$\vb{v}\in\R^m$和$\beta\in\R$的Householder参数。
    
    使$\vb{H}=\vb{I}-\beta\vb{v}\vb{v}^T$是Householder矩阵，且满足$\vb{H}\vb{x}=\norm{\vb{x}}_2\vb{e}_1$。
    \begin{Algoplain}
        \Function{$[v,\beta]=\mathrm{house}(x)$}{
            $\vb{v}(1)=1,\ \vb{v}(2:m)=\vb{x}(2:m),\ \sigma=\vb{x}(2:m)\vb{x}(2:m)^T$\;
            \uIf{$\sigma=0$ and $\vb{x}(1)\geq 0$}{$\beta=0$\;}
            \uElseIf{$\sigma=0$ and $\vb{x}(1)< 0$}{$\beta=-2$\;}
            \Else{
                $\mu=\sqrt{\vb{x}(1)^2+\sigma}$\;
                \uIf{$\vb{x}(1)\leq 0$}{$\vb{v}(1)=\vb{x}(1)-\mu$}   
                \Else{$\vb{v}(1)=-\sigma/(\vb{x}(1)+\mu)$}
                $\beta=2\vb{v}(1)^2/(\sigma+\vb{v}(1)^2)$\;
                $\vb{v}=\vb{v}/\vb{v}(1)$\;
            }
        }
    \end{Algoplain}
\end{BoxAlgorithm}

然而，在算法层面，直接计算$\vb{v}=\vb{x}-\norm{\vb{x}}_2\vb{e}_1$有一些额外的麻烦，注意到
\begin{Equation}[Householder参数1]
    \vb{v}(2:m)=\vb{x}(2:m)
\end{Equation}

因此其实真正要关心的只有$\vb{v}(1)$，我们会分为$\vb{x}(1)\leq 0$和$\vb{x}(1)>0$两种情况来讨论。

对于$\vb{x}(1)\leq 0$，直接进行计算
\begin{Equation}[Householder参数2]
    \vb{v}(1)=\vb{x}(1)-\norm{\vb{x}}_2
\end{Equation}

对于$\vb{x}(1)>0$，考虑进行如下变换以避免数值精度问题
\begin{Equation}[Householder参数3]
    \vb{v}(1)=\frac{\vb{x}(1)-\norm{\vb{x}}_2^2}{\vb{x}(1)+\norm{\vb{x}}_2}=-\frac{\vb{x}(2)^2+\cdots+\vb{x}(m)^2}{\vb{x}(1)+\norm{\vb{x}}_2}
\end{Equation}

引入代换变量$\sigma$和$\mu$，并注意到
\begin{Equation}[Householder参数4]
    \sigma=\vb{x}(2)^2+\cdots+\vb{x}(m)^2\qquad
    \mu=\sqrt{\vb{x}(1)+\sigma}=\norm{\vb{x}}_2
\end{Equation}

这样一来，\cref{eq:Householder参数2,eq:Householder参数3}就可以表示为
\begin{Equation}[Householder参数5]
    \vb{v}(1)=\begin{cases}
        \vb{x}(1)-\mu,&\vb{x}(1)\leq 0\\
        -\sigma/(\vb{x}(1)+\mu),&\vb{x}(1)>0\\
    \end{cases}
\end{Equation}

除此之外，在\cref{alg:Householder参数}中，还有一些值得注意的细节。算法期望缩放$\vb{v}$使$\vb{v}(1)=1$，目的是存储$\vb{v}$时可以少用一位，这固然可行，但这种缩放一定会影响$\vb{H}=\vb{I}-\beta\vb{v}\vb{v}^T$的结果。由于$\beta$也由算法给出，因此在计算$\beta$时额外引入了$\vb{v}(1)^2$以补偿稍后$\vb{v}=\vb{v}/\vb{v}(1)$在$\beta\vb{v}\vb{v}^T$上的损失。

\subsection{Givens旋转变换}

在$\R^m$下，通常旋转仍是对其中特定的两个元素进行，需要将\cref{eq:2x2正交旋转}放到$\R^{m\times m}$中完成。

\begin{BoxFormula}[Givens旋转变换]
    Givens旋转变换矩阵$\vb{G}\in\R^{m\times m}$可以表示为以下形式
    \begin{Equation}[Givens旋转变换]
        \vb{G}(i,j,\theta)=
        \begin{bNiceMatrix}[first-row, last-col]
                    &i      &       &j      &       &       \\
            \vb{I}  &       &       &       &       &       \\
                    &+c     &       &+s     &       &i      \\
                    &       &\vb{I} &       &       &       \\
                    &-s     &       &+c     &       &j      \\
                    &       &       &       &\vb{I} &       \\
        \end{bNiceMatrix}
    \end{Equation}
    其中$\theta$是旋转角（$c=\cos\theta,\ s=\sin\theta$），而$i,j$是将要旋转的元素下标。
\end{BoxFormula}

现在我们想要达成的是，对于任意的$\vb{x}$，通过对$\vb{x}(i)$和$\vb{x}(j)$的旋转使后者归零。
\begin{Equation}[旋转变换1]
    \mqty[
        \vb{x}(1)\\
        \vdots\\
        \vb{x}(i)\\
        \vdots\\
        \vb{x}(j)\\
        \vdots\\
        \vb{x}(m)
    ]\to
    \mqty[
        \vb{x}(1)\\
        \vdots\\
        \times\\
        \vdots\\
        0\\
        \vdots\\
        \vb{x}(m)
    ]
\end{Equation}

请注意Givens矩阵$\vb{G}$表旋转示$\theta$时需要$\vb{G}^T$左乘，不妨记
\begin{Equation}[旋转变换2]
    \vb{G}(i,k,\theta)^T\vb{x}=\vb{y}
\end{Equation}

注意到$\vb{y}(k)$满足
\begin{Equation}[旋转变换3]
    \vb{y}(k)=\begin{cases}
        c\vb{x}(i)-s\vb{x}(j),&k=i\\
        s\vb{x}(i)+c\vb{x}(j),&k=j\\
        \vb{x}(k),&k\neq i,j\\
    \end{cases}
\end{Equation}

通过\cref{eq:旋转变换3}，可以确定，如果要令$\vb{y}(j)=0$，应当将$c,s$取为
\begin{Equation}[旋转变换4]
    c=\frac{\vb{x}(i)}{\sqrt{\vb{x}(i)^2+\vb{x}(j)^2}}\qquad
    s=\frac{-\vb{x}(j)}{\vb{x}(i)^2+\vb{x}(j)^2}
\end{Equation}

我们可以将结论整理如下
\begin{BoxFormula}[Givens矩阵的构造]
    对于$\vb{x}\in\R^m$，若选取以下Givens参数$c,s$
    \begin{Equation}[Givens矩阵的构造1]
        c=\frac{\vb{x}(i)}{\sqrt{\vb{x}(i)^2+\vb{x}(j)^2}}\qquad
        s=\frac{-\vb{x}(j)}{\sqrt{\vb{x}(i)^2+\vb{x}(j)^2}}
    \end{Equation}
    并记$c,s$构造的Givens矩阵作用在$\vb{x}$上的变换结果为$\vb{y}$
    \begin{Equation}[Givens矩阵的构造2]
        \vb{G}^T\vb{x}=\vb{y}
    \end{Equation}
    则其满足
    \begin{Equation}[Givens矩阵的构造3]
        \vb{y}(j)=0
    \end{Equation}
\end{BoxFormula}

在算法层面，只关注需要旋转的元素，且有一些技巧可以减少运算并提高数值精度。

\begin{BoxAlgorithm}[Givens参数]
    对于$\vb{x}\in\R^2$，该函数计算$c\in\R$和$s\in\R$的Givens参数，并满足
    \begin{Equation}[Givens参数算法]
        \mqty[
            c & s \\ 
            -s & c \\
        ]^T
        \mqty[
            a\\
            b\\
        ]=
        \mqty[
            r\\
            0\\
        ]
    \end{Equation}
    \begin{Algoplain}
        \Function{$[c,s]=\mathrm{givens}(x)$}{
            \uIf{$b=0$}{$c=1,\ s=0$}
            \Else{
                \uIf{$\abs{b}>\abs{a}$}{$\tau=-a/b,\ s=1/\sqrt{1+\tau^2},\ c=s\tau$\;}
                \Else{$\tau=-b/a,\ c=1/\sqrt{1+\tau^2},\ s=c\tau$\;}
            }
        }
    \end{Algoplain}
\end{BoxAlgorithm}

\section{QR分解}

QR分解是适用于$m>n$的高矩阵$\vb{A}\in\R^{m\times n}$的分解$\vb{A}=\vb{Q}\vb{R}$，其中，$\vb{Q}\in\R^{m\times m}$是正交矩阵，$\vb{R}\in\R^{m\times n}$是高的上三角矩阵。本质上来说，QR分解其实也是一个将$\vb{A}$通过消元上三角化为$\vb{R}$的过程。消元依赖的就是\cref{sec:两种正交变换}介绍的Householder反射和Givens旋转，它们可以分别将一列或一个元素消为$0$，相乘就构成了$\vb{Q}$。本节先讨论QR分解的存在性和唯一性。

\subsection{QR分解的存在性}

\begin{BoxTheorem}[QR分解]
    若$\vb{A}\in\R^{m\times n}$，则存在正交矩阵$\vb{Q}=\R^{m\times m}$和上三角矩阵$\vb{R}\in\R^{m\times n}$，使
    \begin{Equation}[QR分解]
        \vb{A}=\vb{Q}\vb{R}
    \end{Equation}
\end{BoxTheorem}

\begin{Proof}[\cref{thm:QR分解}]
    使用数学归纳法，先考虑$n=1$即$\vb{A}$只有一列的情况，此时目标的$\vb{R}$实际上就是正比于$\vb{e}_1$的列向量。此时，只要构造Householder矩阵$\vb{Q}^T$使$\vb{R}=\vb{Q}^T\vb{A}$，就满足$\vb{R}(2:m)=\vb{0}$。由于正交矩阵的逆$\vb{Q}^{-1}$和其转置$\vb{Q}^T$相同，这就证明了$n=1$的情形下存在$\vb{A}=\vb{Q}\vb{R}$的分解。

    而对于$n>1$的情况，假设QR分解对于$n$成立，证明其对于$n+1$成立，将$\vb{A}$划分为
    \begin{Equation}[QR分解的存在性1]
        \vb{A}=\qty[
            \begin{array}{c|c}
                \vb{A}_1 & \vb{v} \\
            \end{array}
        ]
    \end{Equation}
    
    其中$\vb{A}\in\R^{m\times (n+1)}$，而$\vb{A}_1=\vb{A}(:,n)\in\R^{m\times n}$，存在QR分解
    \begin{Equation}[QR分解的存在性2]
        \vb{A}_1=\vb{Q}_1\vb{R}_1
    \end{Equation}
    
    或者写为
    \begin{Equation}[QR分解的存在性3]
        \vb{R}_1=\vb{Q}_1^T\vb{A}_1
    \end{Equation}

    当然，我们也可以将$\vb{Q}_1^T$作用在$\vb{v}$上
    \begin{Equation}[QR分解的存在性4]
        \vb{w}=\vb{Q}_1^T\vb{v}
    \end{Equation}

    这里$\vb{w}$显然不会是我们期望的分解，然而，不妨进一步对$\vb{w}(n:m)$做QR分解
    \begin{Equation}[QR分解的存在性5]
        \vb{w}(n:m)=\vb{Q}_2\vb{R}_2
    \end{Equation}

    我们宣称构造这样的$\vb{Q}$和$\vb{R}$，容易验证其分别是正交矩阵和上三角矩阵
    \begin{Equation}[QR分解的存在性6]
        \vb{Q}=\vb{Q}_1\mqty[
            \vb{I}&\vb{0}\\
            \vb{0}&\vb{Q}_2\\
        ]\qquad
        \vb{R}=\qty[
            \begin{array}{c|c}
                \multirow{2}{*}{$\vb{R}_1$} & \vb{w}(1:n-1)\\
                & \vb{R}_2 \\
            \end{array}
        ]
    \end{Equation}

    我们现在验证\cref{eq:QR分解的存在性6}给出的$\vb{Q},\vb{R}$，满足$\vb{A}=\vb{Q}\vb{R}$
    \begin{Equation}[QR分解的存在性7]
        \vb{Q}\vb{R}=\vb{Q}_1\mqty[
            \vb{I}&\vb{0}\\
            \vb{0}&\vb{Q}_2\\
        ]\qty[
            \begin{array}{c|c}
                \multirow{2}{*}{$\vb{R}_1$} & \vb{w}(1:n-1)\\
                & \vb{R}_2 \\
            \end{array}
        ]
    \end{Equation}

    很容易将后两个矩阵相乘
    \begin{Equation}[QR分解的存在性8]
        \vb{Q}\vb{R}=\vb{Q}_1\qty[
            \begin{array}{c|c}
                \multirow{2}{*}{$\vb{R}_1$} & \vb{w}(1:n-1)\\
                & \vb{Q}_2\vb{R}_2 \\
            \end{array}
        ]
    \end{Equation}

    根据\cref{eq:QR分解的存在性5}，这里$\vb{Q}_2\vb{R}_2$实际上就是$\vb{w}(n:m)$
    \begin{Equation}[QR分解的存在性9]
        \vb{Q}\vb{R}=\vb{Q}_1
        \qty[
            \begin{array}{c|c}
                \vb{R}_1 & \vb{w}
            \end{array}
        ]=
        \qty[
            \begin{array}{c|c}
                \vb{Q}_1\vb{R}_1 & \vb{Q}_1\vb{w}
            \end{array}
        ]=
        \qty[
            \begin{array}{c|c}
                \vb{A}_1 & \vb{v}
            \end{array}
        ]=
        \vb{A}
    \end{Equation}

    这就证明了$n>1$的情形下$\vb{A}$存在QR分解！
\end{Proof}

\begin{BoxProperty}[QR分解的性质1]
    若$\vb{A}\in\R^{m\times n}$列满秩，$\vb{A}=\vb{Q}\vb{R}$是QR分解，并记
    \begin{Equation}[QR分解列划分]
        \vb{A}=\qty[
            \begin{array}{c|c|c}
                \vb{a}_1 & \cdots & \vb{a}_n \\
            \end{array}
        ]\qquad
        \vb{Q}=\qty[
            \begin{array}{c|c|c|c|c}
                \vb{q}_1 & \cdots & \vb{q}_n & \cdots & \vb{q}_m \\
            \end{array}
        ]
    \end{Equation}
    那么，对于$j:1:n$，有以下结论成立
    \begin{Gather}[QR分解的性质1]
        \vspan\qty{\vb{a}_1,\cdots,\vb{a}_j}=\vspan\qty{\vb{q}_1,\cdots,\vb{q}_j}\id{a}\\
        r_{jj}\neq 0\id{b}
    \end{Gather}
\end{BoxProperty}

\begin{Proof}[\cref{ppt:QR分解的性质1}]
    不妨根据$\vb{A}=\vb{Q}\vb{R}$考虑$\vb{A}$的第$j$列
    \begin{Equation}[QR分解的性质1-1]
        \vb{a}_j=\Sum[i=1][j]r_{ij}\vb{q}_i\in\vspan\qty{\vb{q}_1,\cdots,\vb{q}_k}
    \end{Equation}

    这就告诉我们
    \begin{Equation}[QR分解的性质1-2]
        \vspan\qty{\vb{a}_1,\cdots,\vb{a}_j}\subseteq\vspan\qty{\vb{q}_1,\cdots,\vb{q}_j}
    \end{Equation}

    而$\vb{A}$作为满秩矩阵，其前$j$个列向量是线性无关的，生成空间的维度为$j$，故有
    \begin{Equation}[QR分解的性质1-3]
        \vspan\qty{\vb{a}_1,\cdots,\vb{a}_j}=\vspan\qty{\vb{q}_1,\cdots,\vb{q}_j}
    \end{Equation}

    同时，在\cref{eq:QR分解的性质1-1}中，若$r_{jj}=0$，则$\vb{a}_1,\cdots,\vb{a}_j$就是线性相关的，违背满秩条件。
\end{Proof}

\begin{BoxProperty}[QR分解的性质2]
    若$\vb{A}\in\R^{m\times n}$列满秩，$\vb{A}=\vb{Q}\vb{R}$是QR分解，并记
    \begin{Equation}[QR分解Q划分]
        \vb{Q}_1=\vb{Q}(:,1:n)\qquad
        \vb{Q}_2=\vb{Q}(:,n+1:m)
    \end{Equation}
    则有
    \begin{Gather}[QR分解的性质2]
        \rank(\vb{A})=\rank(\vb{Q}_1)\id{a}\\
        \rank(\vb{A})^{\perp}=\rank(\vb{Q}_2)\id{b}
    \end{Gather}
\end{BoxProperty}

\begin{Proof}[\cref{ppt:QR分解的性质2}]
    根据\cref{ppt:QR分解的性质1}，有$\vspan\qty{\vb{a}_1,\cdots,\vb{a}_n}=\vspan\qty{\vb{q}_1,\cdots,\vb{q}_n}$，因此$\rank(\vb{A})=\rank(\vb{Q}_1)$。而$\vb{A}$的正交补空间，则是由$\R^{m}$中剩下的$m-n$个正交向量$\vb{Q}_2$表征，因此$\rank(\vb{A})^{\perp}=\rank(\vb{Q}_2)$。
\end{Proof}

\cref{ppt:QR分解的性质2}也给了我们了一个启示，尽管$\vb{A}=\vb{Q}\vb{R}$中$\vb{Q}\in\R^{m\times m}$，但其实表达$\vb{A}$只需要$\vb{Q}$的前$n$列的$\vb{Q}\in\R^{m\times n}$就够了！而QR分解本身也可以仅由$\vb{Q}_1$来表示，这就是Thin QR！

\begin{BoxTheorem}[Thin QR分解]
    若$\vb{A}\in\R^{m\times n}$，则其QR分解还可以写为Thin QR的形式
    \begin{Equation}[Thin QR分解]
        \vb{A}=\vb{Q}\vb{R}=
        \qty[
            \begin{array}{c|c}
                \vb{Q}_1 & \vb{Q}_2\\
            \end{array}
        ]
        \qty[
            \begin{array}{c}
                \vb{R}_1 \\
                \vb{0} \\
            \end{array}
        ]=
        \vb{Q}_1\vb{R}_1
    \end{Equation}
    其中，$\vb{Q}_1\in\R^{m\times n}$，$\vb{Q}_2\in\R^{m\times(m-n)}$，$\vb{R}_1\in\R^{n\times n}$。

    若$\vb{A}$是满秩的，则$\vb{Q}_1,\vb{R}_1$是唯一的，且$\vb{R}_1$主对角线元素为正。
\end{BoxTheorem}

\begin{Proof}[\cref{thm:Thin QR分解}]
    Thin QR分解的形式是显然的，本质上和$\vb{Q}_2$相乘的那部分是作为高的上三角矩阵$\vb{R}$中恒定是零的下半部分，因此$\vb{A}=\vb{Q}\vb{R}=\vb{Q}_1\vb{R}_2+\vb{Q}_2\vb{0}=\vb{Q}_1\vb{R}_1$其实就是滤除了这部分无效计算。

    Thin QR分解的唯一性要借助Cholesky分解来证明，注意到
    \begin{Equation}[Thin QR分解1]
        \vb{A}^T\vb{A}=(\vb{Q}_1\vb{R}_1)^T(\vb{Q}_1\vb{R}_1)=\vb{R}_1^T\vb{Q}_1^T\vb{Q}_1\vb{R}_1=\vb{R}_1^T\vb{R}_1
    \end{Equation}

    根据\cref{thm:Cholesky分解}，$\vb{G}=\vb{R}_1^T$是$\vb{A}^T\vb{A}$的Choleskey分解矩阵，这就论证了$\vb{R}_1$是主对角线元素为正的上三角矩阵并且$\vb{R}_1$唯一。最后，再结合$\vb{Q}_1=\vb{A}\vb{R}_1^{-1}$，这就论证了$\vb{Q}_1$也是唯一的。
\end{Proof}

Thin QR是$\R^{m\times n}=\R^{m\times n}\times\R^{n\times n}$的过程，QR是$\R^{m\times n}=\R^{m\times m}\times\R^{m\times n}$的过程。

\section{基于Householder反射变换的QR分解}
\section{基于Givens旋转变换的QR分解}
\section{基于Gram-Schmidt正交化的QR分解}
\section{最小二乘问题}

\subsection{最小二乘问题}

不妨考虑这样的一个场景
\begin{Equation}[最小二乘1]
    \vb{y}=\vb{A}\vb{x}+\vb{v}
\end{Equation}

设$\vb{A}\in\R^{m\times n}$，当$m>n$即方程数多于未知数时，系统是被过约束的。由于这类情况往往出现在观测数据中，常会伴有噪声$\vb{v}$，作为线性方程组通常是无解的。此时，问题就从精确求解出$\vb{x}$演变了成找到这样一个$\vb{x}$使得$\vb{y}-\vb{A}\vb{x}$最小，这就是最小二乘（Least Square）问题。
\begin{Equation}[最小二乘2]
    \min_{\vb{x}\in\R^n}\norm{\vb{y}-\vb{A}\vb{x}}_2^2
\end{Equation}

\subsection{最小二乘定理}

最小二乘问题虽然表面上是一个优化问题，但可以将其求解转化为精确的线性方程组！

\begin{BoxTheorem}[最小二乘定理]
    若$\vb{x}_{LS}$是最小二乘问题的解
    \begin{Equation}[最小二乘定理A]
        \min_{\vb{x}\in\R^n}\norm{\vb{y}-\vb{A}\vb{x}}_2^2
    \end{Equation}
    当且仅当
    \begin{Equation}[最小二乘定理B]
        \vb{A}^T\vb{A}\vb{x}_{LS}=\vb{A}^T\vb{y}
    \end{Equation}
    特别的，若$\vb{A}$是列满秩的
    \begin{Equation}[最小二乘定理C]
        \vb{x}_{LS}=(\vb{A}^T\vb{A})^{-1}\vb{A}^T\vb{y}
    \end{Equation}
\end{BoxTheorem}

\begin{Proof}[\cref{thm:最小二乘定理}]
    最小二乘问题是指，我们要找这样一个$\vb{x}_{LS}$使$\vb{A}\vb{x}_{LS}$尽可能接近$\vb{y}$。这相当于令$\vb{z}=\vb{A}\vb{x}$，即使$\vb{z}\in\rangespace(\vb{A})$，随后最小化$\vb{z}-\vb{y}$，而这其实就是在寻找$\vb{y}$在$\rangespace(\vb{A})$上的投影！因此有以下关系
    \begin{Equation}[最小二乘定理1]
        \Pi_{\rangespace(\vb{A})}(\vb{y})=\arg\min_{\vb{z}\in\rangespace(\vb{A})}\norm{\vb{z}-\vb{y}}_2^2=\vb{A}\vb{x}_{LS}
    \end{Equation}
    根据\cref{thm:投影定理}的投影定理，对于任意的$\vb{z}\in\rangespace(\vb{A})$
    \begin{Equation}[最小二乘定理2]
        \vb{z}^T(\vb{A}\vb{x}_{LS}-\vb{y})=0
    \end{Equation}
    由于$\vb{z}=\vb{A}\vb{x}$，这等价于对于任意的$\vb{x}\in\R^n$
    \begin{Equation}[最小二乘定理3]
        \vb{x}^T\vb{A}^T(\vb{A}\vb{x}_{LS}-\vb{y})=0
    \end{Equation}
    既然$\vb{x}\in\R^n$，这就说明剩下的部分恒等于$\vb{0}$
    \begin{Equation}[最小二乘定理4]
        \vb{A}^T(\vb{A}\vb{x}_{LS}-\vb{y})=\vb{0}
    \end{Equation}
    移项即得$\vb{A}^T\vb{A}\vb{x}_{LS}=\vb{A}^T\vb{y}$，这就证明了其是$\vb{x}_{LS}$为最小二乘的解的充要条件。
\end{Proof}

\subsection{正交投影算子}

现在回看证明\cref{thm:最小二乘定理}中的过程，注意到$\vb{y}$在$\rangespace(\vb{A})$和$\rangespace(\vb{A})^{\perp}$上的投影可以表示为
\begin{Gather}[正交投影算子的引入]
    \Pi_{\rangespace(\vb{A})}(\vb{y})=\vb{A}\vb{x}_{LS}=(\vb{A}(\vb{A}^T\vb{A})^{-1}\vb{A}^T)\vb{y}\id{1}\\
    \Pi_{\rangespace(\vb{A})^{\perp}}=\vb{y}-\vb{A}\vb{x}_{LS}=(\vb{I}-\vb{A}(\vb{A}^T\vb{A})^{-1}\vb{A}^T)\vb{y}\id{2}
\end{Gather}

换言之，我们找到了将任何向量投影到$\rangespace(\vb{A})$和$\rangespace(\vb{A})^{\perp}$的矩阵！
\begin{BoxDefinition}[正交投影算子]
    矩阵$\vb{A}$的正交投影算子（Orthogonal Projector）定义为
    \begin{Equation}[正交投影算子]
        \vb{P}_{\vb{A}}=\vb{A}(\vb{A}^T\vb{A})^{-1}\vb{A}^{T}\\
    \end{Equation}
    矩阵$\vb{A}$的正交补投影算子（Orthogonal Complement Projector）定义为
    \begin{Equation}[正交补投影算子]
        \vb{P}_{\vb{A}}^{\perp}=\vb{I}-\vb{A}(\vb{A}^T\vb{A})^{-1}\vb{A}^T
    \end{Equation}
\end{BoxDefinition}

投影算子具有幂等性！因为一个向量投影显然等于向量投影的投影。

\begin{BoxProperty}[正交投影算子的幂等性]
    若$\vb{P}_{\vb{A}}$是一个正交投影算子，则其具有幂等性
    \begin{Equation}[正交投影算子的幂等性]
        \vb{P}_{\vb{A}}\vb{P}_{\vb{A}}=\vb{P}_{\vb{A}}
    \end{Equation}
\end{BoxProperty}

\begin{BoxProperty}[正交投影算子的对称性]
    若$\vb{P}_{\vb{A}}$是一个正交投影算子，则其具有对称性
    \begin{Equation}[正交投影算子的对称性]
        \vb{P}_{\vb{A}}^T=\vb{P}_{\vb{A}}
    \end{Equation}
\end{BoxProperty}

\subsection{伪逆}

通常来说，逆矩阵是对于方阵定义的，而对于非方阵，可以引入伪逆的概念。

\begin{BoxDefinition}[伪逆]
    若矩阵$\vb{A}\in\R^{m\times n}$是列满秩的，定义其伪逆（Pseudo Inverse）
    \begin{Equation}[伪逆]
        \vb{A}^{\dagger}=(\vb{A}^T\vb{A})^{-1}\vb{A}^T
    \end{Equation}
\end{BoxDefinition}

伪逆满足部分的逆性，有$\vb{A}^{\dagger}\vb{A}=\vb{I}$但未必有$\vb{A}\vb{A}^{\dagger}=\vb{I}$成立。

\begin{BoxProperty}[伪逆的逆性]
    伪逆满足部分的逆性
    \begin{Equation}[伪逆的逆性]
        \vb{A}^{\dagger}\vb{A}=\vb{I}
    \end{Equation}
\end{BoxProperty}

伪逆可以简化许多结论的表达：最小二乘的解$\vb{x}_{LS}=\vb{A}^{\dagger}\vb{y}$，正交投影算子$\vb{P}_{\vb{A}}=\vb{A}\vb{A}^{\dagger}$。

\subsection{应用LU分解计算最小二乘法}

关于最小二乘法的求解，最直接的方式是对\cref{thm:最小二乘定理}用Cholesky分解
\begin{Equation}[应用LU分解计算最小二乘法]
    \vb{A}^T\vb{A}\vb{x}_{LS}=\vb{A}^T\vb{y}
\end{Equation}

该方法的时间复杂度是$\Obig(mn^2+n^3/3)$，其中，$\mathcal{O}(mn^2)$消耗在计算$\vb{C}=\vb{A}^T\vb{A}$，$\mathcal{O}(n^3/3)$消耗在计算$\vb{C}=\vb{G}\vb{G}^T$的Cholesky分解。没错！计算$\vb{A}^T\vb{A}$和完成Cholesky分解同样耗时！

\subsection{应用QR分解计算最小二乘法}

使用QR分解，从最小二乘法的定义出发，可以更快的计算出结果
\begin{Equation}[应用QR分解计算最小二乘法1]
    \min_{\vb{x}\in\R^n}\norm{\vb{y}-\vb{A}\vb{x}}_2^2
\end{Equation}

若已知$\vb{A}=\vb{Q}\vb{R}$，由于$\vb{Q}$是正交矩阵，不改变范数
\begin{Equation}[应用QR分解计算最小二乘法2]
    \norm{\vb{y}-\vb{A}\vb{x}}_2^2=\norm{\vb{Q}^T\vb{y}-\vb{Q}^T\vb{A}\vb{x}}_2^2
\end{Equation}

现在代入$\vb{A}=\vb{Q}\vb{R}$
\begin{Equation}[应用QR分解计算最小二乘法3]
    \norm{\vb{y}-\vb{A}\vb{x}}_2^2=\norm{\vb{Q}^T\vb{y}-\vb{R}\vb{x}}_2^2
\end{Equation}

根据QR分解和Thin QR分解的关系
\begin{Equation}[应用QR分解计算最小二乘法4]
    \norm{\vb{y}-\vb{A}\vb{x}}_2^2=\norm{
        \mqty[
            \vb{Q}_1^T\vb{y}\\
            \vb{Q}_2^T\vb{y}\\
        ]-
        \mqty[
            \vb{R}_1\vb{x}\\
            \vb{0}
        ]
    }
\end{Equation}

即
\begin{Equation}[应用QR分解计算最小二乘法5]
    \norm{\vb{y}-\vb{A}\vb{x}}_2^2=\norm{\vb{Q}_1^T\vb{y}-\vb{R}_1\vb{x}}+\norm{\vb{Q}_2^T\vb{y}}_2^2
\end{Equation}

因此，最小化$\vb{y}-\vb{A}\vb{x}$的任务就变成了求解$\vb{R}_1\vb{x}=\vb{Q}_1^T\vb{y}$！

该方法的时间复杂度是$\mathcal{O}(2mn^2-2n^2/3)$，这完全来自仅计算$\vb{R}$的Householder QR的复杂度。请注意，计算$\vb{Q}_1^T\vb{y}$不需要显式构造$\vb{Q}$，因为Householder向量已经被存储，利用这些向量计算$\vb{Q}_1^T\vb{y}$的复杂度远低于构造出$\vb{Q}$的复杂度。需要指出的是，QR分解比Cholesky分解在求解最小二乘问题上更慢！对于$m\gg n$的情况，QR分解需要$\mathcal{O}(2mn^2)$，Cholesky分解只需要$\mathcal{O}(mn^2)$。不过，QR分解的优势在于不需要计算$\vb{A}^T\vb{A}$，在数值稳定性上更有优势。
